\documentclass[11pt,a4paper]{article}
\usepackage[utf8]{inputenc}
\usepackage[T1]{fontenc}
\usepackage[catalan]{babel}
\usepackage{lmodern}
\usepackage[sfdefault,lining]{FiraSans}
\usepackage[a4paper,top=2.3cm,bottom=2.2cm,left=1.8cm,right=1.8cm,headheight=1.3cm,headsep=0.4cm,footskip=1cm]{geometry}
\usepackage{amsmath,amssymb}
\usepackage{icomma}
\usepackage{graphicx}
\usepackage[table]{xcolor}
\usepackage{tikz}
\usepackage{fontawesome5}
\usepackage{ccicons}
\usepackage{fancyhdr}
\usepackage{lastpage}
\usepackage{tasks}
\usepackage{needspace}
\usepackage[most]{tcolorbox}
\usepackage{array}
\usepackage[hidelinks,pdftitle={Solucions de la fitxa de la unitat 3: Polinomis. 4t d'ESO Matemàtiques A},pdfsubject={Matemàtiques A 4t ESO},pdfkeywords={polinomis, monomis, valor numèric, suma de polinomis, resta de polinomis, producte de polinomis, identitats notables, divisió de polinomis, regla de Ruffini, teorema del residu, arrels, factorització, factor comú}]{hyperref}

\definecolor{navy}{HTML}{1B2748}
\definecolor{accent}{HTML}{E8603A}
\definecolor{teal}{HTML}{1F9E90}
\definecolor{soft}{HTML}{EEF1F8}
\definecolor{muted}{HTML}{5D6685}

\setlength{\parindent}{0pt}
\setlength{\parskip}{3pt}
\settasks{label=\alph*), label-format=\bfseries\color{navy}, label-width=1.3em, item-indent=1.8em, column-sep=0.6em, after-item-skip=5pt, before-skip=3pt, after-skip=4pt}

% ---------- capçalera i peu ----------
\newcommand{\llicencia}{\href{https://creativecommons.org/licenses/by-nc-sa/4.0/deed.ca}{\ccbyncsa}\ \ Llicència Creative Commons Reconeixement-NoComercial-CompartirIgual 4.0 Internacional (CC BY-NC-SA 4.0)}
\pagestyle{fancy}
\fancyhf{}
\fancyhead[L]{\raisebox{-0.35\height}{\includegraphics[height=0.95cm]{logo-gva.png}}}
\fancyhead[C]{\small\color{muted}Solucions de la fitxa · Unitat 3: Polinomis · 4t ESO A}
\fancyhead[R]{\raisebox{-0.35\height}{\includegraphics[height=0.95cm]{logo-vermellar.png}}}
\fancyfoot[L]{\scriptsize\color{muted}\llicencia}
\fancyfoot[R]{\small\color{muted}\thepage}
\renewcommand{\headrulewidth}{0.4pt}
\renewcommand{\headrule}{\color{muted!40}\hrule height \headrulewidth}
\fancypagestyle{primera}{\fancyhf{}\renewcommand{\headrulewidth}{0pt}
  \fancyfoot[L]{\scriptsize\color{muted}\llicencia}
  \fancyfoot[R]{\small\color{muted}\thepage}}

% ---------- nivells i exercicis ----------
\newcommand{\estrelles}[1]{{\color{accent}\ifcase#1\or\faStar\faStar[regular]\faStar[regular]\or\faStar\faStar\faStar[regular]\or\faStar\faStar\faStar\fi}}
\newcounter{exer}
\newcommand{\ex}[1]{\par\needspace{4\baselineskip}\medskip\refstepcounter{exer}\noindent{\color{navy}\bfseries\theexer.}\ \ {\small\estrelles{#1}}\ \ }

\newcommand{\seccio}[2]{\par\needspace{16\baselineskip}\bigskip
  \begin{tcolorbox}[colback=navy,colframe=navy,coltext=white,arc=2mm,boxrule=0pt,left=3mm,right=3mm,top=1.2mm,bottom=1.2mm]
  {\large\bfseries #1\quad #2}
  \end{tcolorbox}\vspace{-2pt}}

\newcommand{\exf}[1]{\par\medskip\refstepcounter{exer}\noindent{\color{navy}\bfseries\theexer.}\ \ {\small\estrelles{#1}}\ \ }
\newcolumntype{C}[1]{>{\centering\arraybackslash}p{#1}}
% quadre de recordatori (només a la fitxa d'exercicis)
\newcommand{\recorda}[1]{\par\smallskip\begin{tcolorbox}[colback=accent!7,colframe=accent!55,arc=2mm,boxrule=0.5pt,left=2.5mm,right=2.5mm,top=1mm,bottom=1mm]\small{\color{accent}\faLightbulb[regular]}\ \textbf{Recorda:}\ #1\end{tcolorbox}}
\newcommand{\res}[1]{\boxed{#1}}\newcommand{\rest}[1]{\fbox{#1}}

\def\fmneed{11}% que el títol d'una secció no quede sol al final del full
% =====================================================================
% Capçalera FentMates de les fitxes: un color per a cada curs.
% Els colors són els de la pàgina inicial (LEVELS a assets/js/materials.js);
% el logo ve de _eines/logo/pdf/ (python3 _eines/logo/genera.py).
%
% Ús, al final del preàmbul de la unitat:
%   % =====================================================================
% Capçalera FentMates de les fitxes: un color per a cada curs.
% Els colors són els de la pàgina inicial (LEVELS a assets/js/materials.js);
% el logo ve de _eines/logo/pdf/ (python3 _eines/logo/genera.py).
%
% Ús, al final del preàmbul de la unitat:
%   % =====================================================================
% Capçalera FentMates de les fitxes: un color per a cada curs.
% Els colors són els de la pàgina inicial (LEVELS a assets/js/materials.js);
% el logo ve de _eines/logo/pdf/ (python3 _eines/logo/genera.py).
%
% Ús, al final del preàmbul de la unitat:
%   \input{../comu/fentmates.tex}
%   \fmcurs{1eso}
% i, al començament del document:
%   \capcalera{Matemàtiques}{1r ESO}{Fitxa d'exercicis · Unitat 2}{Divisibilitat}
% Opcional abans de \input: \def\fmneed{14} (espai mínim abans de cada secció).
% =====================================================================
\makeatletter
\def\fm@logodir{../../logo/pdf}
\providecommand{\fmneed}{10}
\def\fm@cursnom{}

% clau del curs, color i color del text damunt del color
\newcommand{\fm@curs}[3]{\@namedef{fm@color@#1}{#2}\@namedef{fm@text@#1}{#3}}
\fm@curs{1eso}{E6B84A}{navy}
\fm@curs{2eso}{3B82F6}{white}
\fm@curs{3eso}{FF7A50}{white}
\fm@curs{4eso-a}{A78BFA}{white}
\fm@curs{4eso-b}{6D28D9}{white}
\fm@curs{1batx-mi}{9B1C1C}{white}
\fm@curs{1batx-ccss}{EF6B6B}{white}
\fm@curs{2batx-mii}{1E6B47}{white}
\fm@curs{2batx-ccss}{2BB3A3}{white}

\newcommand{\fmcurs}[1]{%
  \@ifundefined{fm@color@#1}{\PackageError{fentmates}{Curs desconegut: #1}{}}{}%
  \edef\fm@tmp{\noexpand\definecolor{curs}{HTML}{\@nameuse{fm@color@#1}}%
               \noexpand\colorlet{curstext}{\@nameuse{fm@text@#1}}}\fm@tmp
  \colorlet{cursfosc}{curs!72!black}%
  \def\fm@logo{\fm@logodir/vertical-#1.pdf}%
  \def\fm@marca{\fm@logodir/marca-#1.pdf}}

% ---------- capçalera de la primera pàgina ----------
% #1 matèria (petit), #2 curs, #3 tipus de document i unitat, #4 títol
\newcommand{\capcalera}[4]{%
  \gdef\fm@cursnom{#2}%
  \par\noindent
  \begin{tcolorbox}[colback=curs,colframe=curs,boxrule=0pt,arc=2.2mm,width=3.5cm,height=2.15cm,
      valign=center,left=3mm,right=2mm,top=1mm,bottom=1mm,nobeforeafter,box align=center]
    \color{curstext}{\footnotesize #1}\par\vspace{1pt}{\fontsize{19}{21}\selectfont\bfseries #2\par}%
  \end{tcolorbox}\hspace{5mm}%
  \begin{minipage}[c]{\dimexpr\linewidth-3.5cm-5mm-2.9cm\relax}
    {\color{cursfosc}\small\bfseries\MakeUppercase{#3}}\par\vspace{2pt}
    {\color{navy}\fontsize{20}{23}\selectfont\bfseries #4\par}
  \end{minipage}\hfill
  \begin{minipage}[c]{2.6cm}\raggedleft\includegraphics[width=2.6cm]{\fm@logo}\end{minipage}\par
  \vspace{2pt}}

% ---------- barres de secció del color del curs ----------
\renewcommand{\seccio}[2]{\par\needspace{\fmneed\baselineskip}\bigskip
  \begin{tcolorbox}[colback=curs,colframe=curs,coltext=curstext,arc=2mm,boxrule=0pt,left=3mm,right=3mm,top=1.2mm,bottom=1.2mm]
  {\large\bfseries #1\quad #2}
  \end{tcolorbox}\vspace{-2pt}}

% ---------- estrelles de dificultat del color del curs (a Batxillerat no se'n posen) ----------
\@ifundefined{estrelles}{}{%
  \renewcommand{\estrelles}[1]{{\color{curs}\ifcase#1\or\faStar\faStar[regular]\faStar[regular]\or\faStar\faStar\faStar[regular]\or\faStar\faStar\faStar\fi}}}

% ---------- capçalera i peu de les pàgines ----------
% primera pàgina: logos institucionals al peu
\fancypagestyle{primera}{\fancyhf{}\renewcommand{\headrulewidth}{0pt}
  \fancyfoot[L]{\raisebox{-0.32\height}{\includegraphics[height=0.78cm]{logo-gva.png}}\hspace{5mm}%
    \raisebox{-0.32\height}{\includegraphics[height=0.78cm]{logo-vermellar.png}}}
  \fancyfoot[R]{\scriptsize\color{muted}\href{https://creativecommons.org/licenses/by-nc-sa/4.0/deed.ca}{\ccbyncsa}\ \ CC BY-NC-SA 4.0 · fentmates.org\qquad\small\thepage}}
% la resta: el símbol a l'esquerra i el curs a la dreta, amb el títol al centre com abans
\fancyhead[L]{\raisebox{-0.3\height}{\includegraphics[height=0.62cm]{\fm@marca}}}
\fancyhead[R]{\small\bfseries\color{cursfosc}\fm@cursnom}
\renewcommand{\headrule}{\color{curs!45}\hrule height \headrulewidth}
\makeatother

%   \fmcurs{1eso}
% i, al començament del document:
%   \capcalera{Matemàtiques}{1r ESO}{Fitxa d'exercicis · Unitat 2}{Divisibilitat}
% Opcional abans de \input: \def\fmneed{14} (espai mínim abans de cada secció).
% =====================================================================
\makeatletter
\def\fm@logodir{../../logo/pdf}
\providecommand{\fmneed}{10}
\def\fm@cursnom{}

% clau del curs, color i color del text damunt del color
\newcommand{\fm@curs}[3]{\@namedef{fm@color@#1}{#2}\@namedef{fm@text@#1}{#3}}
\fm@curs{1eso}{E6B84A}{navy}
\fm@curs{2eso}{3B82F6}{white}
\fm@curs{3eso}{FF7A50}{white}
\fm@curs{4eso-a}{A78BFA}{white}
\fm@curs{4eso-b}{6D28D9}{white}
\fm@curs{1batx-mi}{9B1C1C}{white}
\fm@curs{1batx-ccss}{EF6B6B}{white}
\fm@curs{2batx-mii}{1E6B47}{white}
\fm@curs{2batx-ccss}{2BB3A3}{white}

\newcommand{\fmcurs}[1]{%
  \@ifundefined{fm@color@#1}{\PackageError{fentmates}{Curs desconegut: #1}{}}{}%
  \edef\fm@tmp{\noexpand\definecolor{curs}{HTML}{\@nameuse{fm@color@#1}}%
               \noexpand\colorlet{curstext}{\@nameuse{fm@text@#1}}}\fm@tmp
  \colorlet{cursfosc}{curs!72!black}%
  \def\fm@logo{\fm@logodir/vertical-#1.pdf}%
  \def\fm@marca{\fm@logodir/marca-#1.pdf}}

% ---------- capçalera de la primera pàgina ----------
% #1 matèria (petit), #2 curs, #3 tipus de document i unitat, #4 títol
\newcommand{\capcalera}[4]{%
  \gdef\fm@cursnom{#2}%
  \par\noindent
  \begin{tcolorbox}[colback=curs,colframe=curs,boxrule=0pt,arc=2.2mm,width=3.5cm,height=2.15cm,
      valign=center,left=3mm,right=2mm,top=1mm,bottom=1mm,nobeforeafter,box align=center]
    \color{curstext}{\footnotesize #1}\par\vspace{1pt}{\fontsize{19}{21}\selectfont\bfseries #2\par}%
  \end{tcolorbox}\hspace{5mm}%
  \begin{minipage}[c]{\dimexpr\linewidth-3.5cm-5mm-2.9cm\relax}
    {\color{cursfosc}\small\bfseries\MakeUppercase{#3}}\par\vspace{2pt}
    {\color{navy}\fontsize{20}{23}\selectfont\bfseries #4\par}
  \end{minipage}\hfill
  \begin{minipage}[c]{2.6cm}\raggedleft\includegraphics[width=2.6cm]{\fm@logo}\end{minipage}\par
  \vspace{2pt}}

% ---------- barres de secció del color del curs ----------
\renewcommand{\seccio}[2]{\par\needspace{\fmneed\baselineskip}\bigskip
  \begin{tcolorbox}[colback=curs,colframe=curs,coltext=curstext,arc=2mm,boxrule=0pt,left=3mm,right=3mm,top=1.2mm,bottom=1.2mm]
  {\large\bfseries #1\quad #2}
  \end{tcolorbox}\vspace{-2pt}}

% ---------- estrelles de dificultat del color del curs (a Batxillerat no se'n posen) ----------
\@ifundefined{estrelles}{}{%
  \renewcommand{\estrelles}[1]{{\color{curs}\ifcase#1\or\faStar\faStar[regular]\faStar[regular]\or\faStar\faStar\faStar[regular]\or\faStar\faStar\faStar\fi}}}

% ---------- capçalera i peu de les pàgines ----------
% primera pàgina: logos institucionals al peu
\fancypagestyle{primera}{\fancyhf{}\renewcommand{\headrulewidth}{0pt}
  \fancyfoot[L]{\raisebox{-0.32\height}{\includegraphics[height=0.78cm]{logo-gva.png}}\hspace{5mm}%
    \raisebox{-0.32\height}{\includegraphics[height=0.78cm]{logo-vermellar.png}}}
  \fancyfoot[R]{\scriptsize\color{muted}\href{https://creativecommons.org/licenses/by-nc-sa/4.0/deed.ca}{\ccbyncsa}\ \ CC BY-NC-SA 4.0 · fentmates.org\qquad\small\thepage}}
% la resta: el símbol a l'esquerra i el curs a la dreta, amb el títol al centre com abans
\fancyhead[L]{\raisebox{-0.3\height}{\includegraphics[height=0.62cm]{\fm@marca}}}
\fancyhead[R]{\small\bfseries\color{cursfosc}\fm@cursnom}
\renewcommand{\headrule}{\color{curs!45}\hrule height \headrulewidth}
\makeatother

%   \fmcurs{1eso}
% i, al començament del document:
%   \capcalera{Matemàtiques}{1r ESO}{Fitxa d'exercicis · Unitat 2}{Divisibilitat}
% Opcional abans de \input: \def\fmneed{14} (espai mínim abans de cada secció).
% =====================================================================
\makeatletter
\def\fm@logodir{../../logo/pdf}
\providecommand{\fmneed}{10}
\def\fm@cursnom{}

% clau del curs, color i color del text damunt del color
\newcommand{\fm@curs}[3]{\@namedef{fm@color@#1}{#2}\@namedef{fm@text@#1}{#3}}
\fm@curs{1eso}{E6B84A}{navy}
\fm@curs{2eso}{3B82F6}{white}
\fm@curs{3eso}{FF7A50}{white}
\fm@curs{4eso-a}{A78BFA}{white}
\fm@curs{4eso-b}{6D28D9}{white}
\fm@curs{1batx-mi}{9B1C1C}{white}
\fm@curs{1batx-ccss}{EF6B6B}{white}
\fm@curs{2batx-mii}{1E6B47}{white}
\fm@curs{2batx-ccss}{2BB3A3}{white}

\newcommand{\fmcurs}[1]{%
  \@ifundefined{fm@color@#1}{\PackageError{fentmates}{Curs desconegut: #1}{}}{}%
  \edef\fm@tmp{\noexpand\definecolor{curs}{HTML}{\@nameuse{fm@color@#1}}%
               \noexpand\colorlet{curstext}{\@nameuse{fm@text@#1}}}\fm@tmp
  \colorlet{cursfosc}{curs!72!black}%
  \def\fm@logo{\fm@logodir/vertical-#1.pdf}%
  \def\fm@marca{\fm@logodir/marca-#1.pdf}}

% ---------- capçalera de la primera pàgina ----------
% #1 matèria (petit), #2 curs, #3 tipus de document i unitat, #4 títol
\newcommand{\capcalera}[4]{%
  \gdef\fm@cursnom{#2}%
  \par\noindent
  \begin{tcolorbox}[colback=curs,colframe=curs,boxrule=0pt,arc=2.2mm,width=3.5cm,height=2.15cm,
      valign=center,left=3mm,right=2mm,top=1mm,bottom=1mm,nobeforeafter,box align=center]
    \color{curstext}{\footnotesize #1}\par\vspace{1pt}{\fontsize{19}{21}\selectfont\bfseries #2\par}%
  \end{tcolorbox}\hspace{5mm}%
  \begin{minipage}[c]{\dimexpr\linewidth-3.5cm-5mm-2.9cm\relax}
    {\color{cursfosc}\small\bfseries\MakeUppercase{#3}}\par\vspace{2pt}
    {\color{navy}\fontsize{20}{23}\selectfont\bfseries #4\par}
  \end{minipage}\hfill
  \begin{minipage}[c]{2.6cm}\raggedleft\includegraphics[width=2.6cm]{\fm@logo}\end{minipage}\par
  \vspace{2pt}}

% ---------- barres de secció del color del curs ----------
\renewcommand{\seccio}[2]{\par\needspace{\fmneed\baselineskip}\bigskip
  \begin{tcolorbox}[colback=curs,colframe=curs,coltext=curstext,arc=2mm,boxrule=0pt,left=3mm,right=3mm,top=1.2mm,bottom=1.2mm]
  {\large\bfseries #1\quad #2}
  \end{tcolorbox}\vspace{-2pt}}

% ---------- estrelles de dificultat del color del curs (a Batxillerat no se'n posen) ----------
\@ifundefined{estrelles}{}{%
  \renewcommand{\estrelles}[1]{{\color{curs}\ifcase#1\or\faStar\faStar[regular]\faStar[regular]\or\faStar\faStar\faStar[regular]\or\faStar\faStar\faStar\fi}}}

% ---------- capçalera i peu de les pàgines ----------
% primera pàgina: logos institucionals al peu
\fancypagestyle{primera}{\fancyhf{}\renewcommand{\headrulewidth}{0pt}
  \fancyfoot[L]{\raisebox{-0.32\height}{\includegraphics[height=0.78cm]{logo-gva.png}}\hspace{5mm}%
    \raisebox{-0.32\height}{\includegraphics[height=0.78cm]{logo-vermellar.png}}}
  \fancyfoot[R]{\scriptsize\color{muted}\href{https://creativecommons.org/licenses/by-nc-sa/4.0/deed.ca}{\ccbyncsa}\ \ CC BY-NC-SA 4.0 · fentmates.org\qquad\small\thepage}}
% la resta: el símbol a l'esquerra i el curs a la dreta, amb el títol al centre com abans
\fancyhead[L]{\raisebox{-0.3\height}{\includegraphics[height=0.62cm]{\fm@marca}}}
\fancyhead[R]{\small\bfseries\color{cursfosc}\fm@cursnom}
\renewcommand{\headrule}{\color{curs!45}\hrule height \headrulewidth}
\makeatother

\fmcurs{4eso-a}

\begin{document}
\thispagestyle{primera}

\capcalera{Matemàtiques A}{4t ESO}{Solucions de la fitxa · Unitat 3}{\raggedright Polinomis}
\vspace{8pt}

\begin{tcolorbox}[colback=soft,colframe=soft,arc=2mm,left=3mm,right=3mm,top=1.5mm,bottom=1.5mm]
\small Intenta resoldre cada exercici abans de mirar-ne la solució. Si no t'ix el mateix resultat, repassa els passos o consulta els apunts interactius de la unitat.
\end{tcolorbox}

% =====================================================
\seccio{1}{Monomis}

\exf{1} \begin{center}\renewcommand{\arraystretch}{1.5}\begin{tabular}{|C{2.4cm}|C{2.6cm}|C{2.6cm}|C{1.6cm}|}\hline\rowcolor{soft}\textbf{Monomi} & \textbf{Coeficient} & \textbf{Part literal} & \textbf{Grau}\\\hline$-3x^4$ & $\res{-3}$ & $\res{x^4}$ & $\res{4}$\\\hline$x^2y$ & $\res{1}$ & $\res{x^2y}$ & $\res{3}$\\\hline$\frac25x$ & $\res{\frac25}$ & $\res{x}$ & $\res{1}$\\\hline$-a^3b^2$ & $\res{-1}$ & $\res{a^3b^2}$ & $\res{5}$\\\hline$8$ & $\res{8}$ & \rest{no en té} & $\res{0}$\\\hline$-x$ & $\res{-1}$ & $\res{x}$ & $\res{1}$\\\hline\end{tabular}\end{center}

\ex{1} $\res{\left\{3x^2,\ 7x^2,\ \tfrac12x^2\right\}}$, \ $\res{\left\{-5x,\ 2x\right\}}$, \ $\res{\left\{x^3,\ -x^3\right\}}$ (la mateixa part literal).

\ex{1} \begin{tasks}(3)
  \task $\res{7x}$
  \task $\res{-2x^2}$
  \task $\res{x^3+8x}$
  \task $\res{2x^2+2x-3}$
  \task $\res{\frac54x}$
  \task $\res{5xy}$
\end{tasks}

\ex{1} \begin{tasks}(3)
  \task $\res{12x^5}$
  \task $\res{-10x^5}$
  \task $\res{6x^5}$
  \task $\res{\frac32x^3}$
  \task $\res{12x^3y^4}$
  \task $\res{-6x^4}$
\end{tasks}

\ex{2} \begin{tasks}(3)
  \task $\res{4x^3}$
  \task $\res{-5x^5}$
  \task $\res{-3}$
  \task $\res{4x^6}$
  \task $\res{-27x^6}$
  \task $\res{\frac1{16}x^4}$
\end{tasks}

\ex{2} \begin{tasks}(1)
  \task $-2\cdot (-1)^{3}=\res{2};\ \ -2\cdot 2^{3}=\res{-16};\ \ -2\cdot (\frac{1}{2})^{3}=\res{-\frac{1}{4}}$
  \task $\frac{3}{4}\cdot (-2)^{2}=\res{3};\ \ \frac{3}{4}\cdot 0^{2}=\res{0};\ \ \frac{3}{4}\cdot 4^{2}=\res{12}$
\end{tasks}


% =====================================================
\seccio{2}{Polinomis i valor numèric}

\exf{1} \begin{center}\renewcommand{\arraystretch}{1.5}\small\begin{tabular}{|l|C{3.7cm}|C{1.1cm}|C{1.5cm}|C{1.6cm}|C{1.95cm}|}\hline\rowcolor{soft}\textbf{Polinomi} & \textbf{Ordenat} & \textbf{Grau} & \textbf{Terme indep.} & \textbf{Coef. princ.} & \textbf{Complet?}\\\hline$P(x)=3x^2-5x^4+x-7$ & $\res{-5x^4+3x^2+x-7}$ & $\res{4}$ & $\res{-7}$ & $\res{-5}$ & \rest{No}\\\hline$Q(x)=6-2x+x^3$ & $\res{x^3-2x+6}$ & $\res{3}$ & $\res{6}$ & $\res{1}$ & \rest{No}\\\hline$R(x)=4x-x^2+1$ & $\res{-x^2+4x+1}$ & $\res{2}$ & $\res{1}$ & $\res{-1}$ & \rest{Sí}\\\hline$S(x)=2x^5+x^2-3x$ & $\res{2x^5+x^2-3x}$ & $\res{5}$ & $\res{0}$ & $\res{2}$ & \rest{No}\\\hline$U(x)=\frac12x^3-x^2+3x-2$ & $\res{\frac12x^3-x^2+3x-2}$ & $\res{3}$ & $\res{-2}$ & $\res{\frac12}$ & \rest{Sí}\\\hline\end{tabular}\end{center}

\ex{1} Per exemple:\begin{tasks}(2)
  \task $\res{x^3+2x^2-x-2}$
  \task $\res{x^4-7}$
  \task $\res{-3x^2+x+5}$
  \task $\res{x^5-4x^2+x}$
\end{tasks}

\ex{1} \begin{tasks}(1)
  \task $P(0)=0^{3}-2\cdot 0^{2}+4\cdot 0-1=\res{-1}$
  \task $P(1)=1^{3}-2\cdot 1^{2}+4\cdot 1-1=\res{2}$
  \task $P(-1)=(-1)^{3}-2\cdot (-1)^{2}+4\cdot (-1)-1=\res{-8}$
  \task $P(2)=2^{3}-2\cdot 2^{2}+4\cdot 2-1=\res{7}$
  \task $P(-2)=(-2)^{3}-2\cdot (-2)^{2}+4\cdot (-2)-1=\res{-25}$
\end{tasks}

\ex{1} \begin{tasks}(3)
  \task $\res{6},\ \ \res{10},\ \ \res{3}$
  \task $\res{-15},\ \ \res{-3},\ \ \res{0}$
  \task $\res{56},\ \ \res{0},\ \ \res{2}$
\end{tasks}

\ex{2} $P(2)=16-2k+5=21-2k=9\Rightarrow\res{k=6}$.

\ex{2} $P(-1)=1-a-4=-3-a=-6\Rightarrow\res{a=3}$.

\ex{2} $C(10)=2+30+50=\res{82\ \text{€}}$; \ $C(100)=200+300+50=\res{550\ \text{€}}$.

\ex{3} $h(0)=\res{1}$, \ $h(1)=-5+20+1=\res{16}$, \ $h(2)=-20+40+1=\res{21}$, \ $h(4)=-80+80+1=\res{1}$ m.\\[2pt]$h(0)=1$: la pilota es llança des d'1 m d'altura; $h(4)=1$: als 4 s torna a estar a 1 m (l'altura màxima, 21 m, és als 2 s).


% =====================================================
\seccio{3}{Suma i resta de polinomis}

\exf{1} \begin{tasks}(2)
  \task $\res{4x^{3}+2x^{2}+x-11}$
  \task $\res{2x^{3}-6x^{2}+x+1}$
  \task $\res{x^{3}+6x^{2}-3x-7}$
  \task $\res{4x^{3}+4x-10}$
\end{tasks}

\ex{2} \begin{tasks}(1)
  \task $\res{2x^{3}-4x^{2}-2x}$
  \task $\res{6x^{3}+2x^{2}-7x-13}$
  \task $\res{x^{3}-12x^{2}+4x+8}$
\end{tasks}

\ex{2} $Q(x)=\left(3x^3+x^2-5\right)-\left(x^3-2x+4\right)=\res{2x^{3}+x^{2}+2x-9}$.

\ex{2} $P(x)=\res{x^{2}+4x+2}$; \ $P(4)=16+16+2=\res{34\ \text{cm}}$.


% =====================================================
\seccio{4}{Producte de polinomis}

\exf{1} \begin{tasks}(2)
  \task $\res{3x^{3}-6x^{2}+15x}$
  \task $\res{-2x^{5}+8x^{3}-2x^{2}}$
  \task $\res{2x^{3}-3x^{2}+x}$
  \task $\res{-5x^{5}+5x^{4}-15x^{3}}$
  \task $\res{-4x^{3}+12x^{2}-4x}$
  \task $\res{-2x^{5}+x^{3}-7x}$
\end{tasks}

\ex{1} \begin{tasks}(2)
  \task $\res{x^{2}-3x-10}$
  \task $\res{6x^{2}+5x-4}$
  \task $\res{x^{3}-3x^{2}+x-3}$
  \task $\res{x^{3}-x^{2}+x+3}$
  \task $\res{2x^{3}-5x^{2}+6x-8}$
  \task $\res{x^{4}+3x-2}$
\end{tasks}

\ex{2} \begin{tasks}(1)
  \task $\res{2x^{3}-5x^{2}+x+2}$
  \task $\res{x^{3}-7x^{2}+12x-9}$
  \task $\res{3x^{3}-12x^{2}+15x-6}$
\end{tasks}

\ex{2} \begin{tasks}(3)
  \task $\res{x^{2}-7x}$
  \task $\res{10}$
  \task $\res{x^{2}+x}$
\end{tasks}

\ex{3} Àrea: $A(x)=(2x+3)(x-1)=\res{2x^{2}+x-3}$; \ $A(5)=13\cdot4=\res{52\ \text{cm}^2}$.\\[2pt]Volum: $V(x)=\res{x^{3}+7x^{2}+10x}$; \ $V(3)=3\cdot5\cdot8=\res{120\ \text{cm}^3}$.


% =====================================================
\seccio{5}{Identitats notables}

\exf{1} \begin{tasks}(3)
  \task $\res{x^{2}+6x+9}$
  \task $\res{x^{2}-8x+16}$
  \task $\res{x^{2}-25}$
  \task $\res{4x^{2}+4x+1}$
  \task $\res{9x^{2}-12x+4}$
  \task $\res{4x^{2}-49}$
  \task $\res{x^{4}+6x^{2}+9}$
  \task $\res{x^{2}-10x+25}$
  \task $\res{\frac{1}{4}x^{2}+x+1}$
\end{tasks}

\ex{2} \begin{tasks}(2)
  \task $\res{4x}$
  \task $\res{6x+18}$
  \task $\res{5x^{2}-4x-3}$
  \task $\res{x^{2}-14x+17}$
\end{tasks}

\ex{2} \begin{tasks}(1)
  \task $(x+\res{6})^2=x^2+\res{12x}+36$
  \task $(\res{2x}-5)^2=4x^2-\res{20x}+25$
  \task $(3x+\res{4})(3x-\res{4})=\res{9x^2}-16$
\end{tasks}

\ex{3} $(100+1)^2=10000+200+1=\res{10201}$; \ $(100-1)^2=10000-200+1=\res{9801}$; \ $(50+1)(50-1)=2500-1=\res{2499}$.


% =====================================================
\seccio{6}{Divisió de polinomis}

\exf{1} \begin{tasks}(3)
  \task $\res{2x^{2}-3x+1}$
  \task $\res{2x^{3}-3x+1}$
  \task $\res{-4x^{2}+2x-1}$
\end{tasks}

\ex{2} \begin{tasks}(1)
  \task $c(x)=\res{2x+3},\ \ r(x)=\res{-3x+1}$
  \task $c(x)=\res{x^{2}-x-1},\ \ r(x)=\res{2x-1}$
  \task $c(x)=\res{3x^{2}+x+1},\ \ r(x)=\res{6}$
  \task $c(x)=\res{x^{2}-1},\ \ r(x)=\res{0}$
  \task $c(x)=\res{3x^{2}+6x+11},\ \ r(x)=\res{15x-11}$
\end{tasks}

\ex{2} $\left(x^2+1\right)(2x+3)+(-3x+1)=2x^{3}+3x^{2}+2x+3+(-3x+1)=\res{2x^{3}+3x^{2}-x+4}$\ \ \checkmark

\ex{3} $D(x)=\left(x^2-2\right)(x+3)+2x-1=x^3+3x^2-2x-6+2x-1=\res{x^{3}+3x^{2}-7}$.

\ex{3} Altura: $\left(2x^2+7x+3\right):(x+3)=\res{2x+1}$ \ (residu 0); \ perímetre: $2(x+3)+2(2x+1)=\res{6x+8}$.


% =====================================================
\seccio{7}{Regla de Ruffini}

\exf{1} \begin{tasks}[after-item-skip=8pt](2)
  \task $\begin{array}{r|rrr|r} & 1 & -4 & 1 & 6\\ 2 & & 2 & -4 & -6\\ \hline & 1 & -2 & -3 & 0\end{array}$\\[2pt]$c(x)=\res{x^{2}-2x-3},\ \ r=\res{0}$
  \task $\begin{array}{r|rrr|r} & 1 & 2 & -1 & 3\\ -1 & & -1 & -1 & 2\\ \hline & 1 & 1 & -2 & 5\end{array}$\\[2pt]$c(x)=\res{x^{2}+x-2},\ \ r=\res{5}$
  \task $\begin{array}{r|rrr|r} & 2 & -3 & 0 & 4\\ 1 & & 2 & -1 & -1\\ \hline & 2 & -1 & -1 & 3\end{array}$\\[2pt]$c(x)=\res{2x^{2}-x-1},\ \ r=\res{3}$
  \task $\begin{array}{r|rrrr|r} & 1 & 0 & -3 & 2 & -5\\ -2 & & -2 & 4 & -2 & 0\\ \hline & 1 & -2 & 1 & 0 & -5\end{array}$\\[2pt]$c(x)=\res{x^{3}-2x^{2}+x},\ \ r=\res{-5}$
  \task $\begin{array}{r|rrrr|r} & 1 & 0 & 0 & 0 & -16\\ 2 & & 2 & 4 & 8 & 16\\ \hline & 1 & 2 & 4 & 8 & 0\end{array}$\\[2pt]$c(x)=\res{x^{3}+2x^{2}+4x+8},\ \ r=\res{0}$
  \task $\begin{array}{r|rrr|r} & 3 & 0 & -5 & 1\\ -3 & & -9 & 27 & -66\\ \hline & 3 & -9 & 22 & -65\end{array}$\\[2pt]$c(x)=\res{3x^{2}-9x+22},\ \ r=\res{-65}$
\end{tasks}

\ex{2} \begin{tasks}(2)
  \task $\begin{array}{r|rrrrr|r} & 1 & 0 & 0 & 0 & 0 & -1\\ 1 & & 1 & 1 & 1 & 1 & 1\\ \hline & 1 & 1 & 1 & 1 & 1 & 0\end{array}$\\[2pt]$c(x)=\res{x^{4}+x^{3}+x^{2}+x+1},\ \ r=\res{0}$
  \task $\begin{array}{r|rrr|r} & 1 & 0 & 0 & 27\\ -3 & & -3 & 9 & -27\\ \hline & 1 & -3 & 9 & 0\end{array}$\\[2pt]$c(x)=\res{x^{2}-3x+9},\ \ r=\res{0}$
\end{tasks}

\ex{2} \begin{tasks}(2)
  \task $\begin{array}{r|rrr|r} & 1 & -2 & -5 & 6\\ 3 & & 3 & 3 & -6\\ \hline & 1 & 1 & -2 & 0\end{array}$\\[2pt]Residu 0: $\res{\text{sí que és divisible}}$
  \task $\begin{array}{r|rrr|r} & 1 & -2 & -5 & 6\\ -1 & & -1 & 3 & 2\\ \hline & 1 & -3 & -2 & 8\end{array}$\\[2pt]Residu 8: $\res{\text{no és divisible}}$
\end{tasks}

\ex{3} Amb Ruffini, el residu és $8+4k-8+12=4k+12$. Ha de ser 0: $4k+12=0\Rightarrow\res{k=-3}$.


% =====================================================
\seccio{8}{Valor numèric amb Ruffini. Teorema del residu}

\exf{1} \begin{tasks}(1)
  \task $r=P(2)=2^{3}-5\cdot 2+2=\res{0}$
  \task $r=P(-1)=(-1)^{4}+2\cdot (-1)^{3}-(-1)=\res{0}$
  \task $r=P(1)=2\cdot 1^{5}-3\cdot 1^{2}+1=\res{0}$
  \task $r=P(-2)=(-2)^{3}+(-2)^{2}+(-2)+1=\res{-5}$
\end{tasks}

\ex{2} \begin{tasks}(2)
  \task $\begin{array}{r|rrrr|r} & 1 & -2 & -5 & 1 & 3\\ 3 & & 3 & 3 & -6 & -15\\ \hline & 1 & 1 & -2 & -5 & -12\end{array}$\\[2pt]$P(3)=\res{-12}$\\[1pt]{\small $3^{4}-2\cdot 3^{3}-5\cdot 3^{2}+3+3=-12$}
  \task $\begin{array}{r|rrrr|r} & 1 & -2 & -5 & 1 & 3\\ -2 & & -2 & 8 & -6 & 10\\ \hline & 1 & -4 & 3 & -5 & 13\end{array}$\\[2pt]$P(-2)=\res{13}$\\[1pt]{\small $(-2)^{4}-2\cdot (-2)^{3}-5\cdot (-2)^{2}+(-2)+3=13$}
\end{tasks}

\ex{2} $P(-2)=-8+2m+4=2m-4=10\Rightarrow\res{m=7}$.

\ex{3} $P(1)=2+a-1+3=4+a=5\Rightarrow\res{a=1}$.


% =====================================================
\seccio{9}{Arrels d'un polinomi}

\exf{1} $P(1)=1-2-1+2=0$, \ $P(-1)=-1-2+1+2=0$, \ $P(2)=8-8-2+2=0$: \ $\res{\text{les tres són arrels}}$.

\ex{1} $P(0)=0$, \ $P(2)=8+4-12=0$, \ $P(-2)=-8+4+12=8$, \ $P(-3)=-27+9+18=0$: \ $\res{\text{són arrels } 0,\ 2\text{ i }-3}$.

\ex{2} \begin{tasks}(2)
  \task $\res{x=1,\ x=2,\ x=3}$
  \task $\res{x=-2,\ x=-1,\ x=1}$
  \task $\res{x=-2,\ x=1\ (\text{doble})}$
  \task $\res{x=-2,\ x=-1,\ x=1,\ x=2}$
\end{tasks}

\ex{2} \begin{tasks}(2)
  \task $\res{x=3,\ x=4}$
  \task $\res{x=-\frac{1}{2},\ x=1}$
  \task $\res{\text{no en té } (\Delta=4-20<0)}$
  \task $\res{x=3\ (\text{doble})}$
\end{tasks}

\ex{2} Les arrels enteres han de ser divisors del terme independent, 6, i \\[2pt]$\res{5\text{ no és divisor de }6}$.

\ex{3} $(x-2)(x+1)(x-3)=\res{x^{3}-4x^{2}+x+6}$; \ $8+2k-6=0\Rightarrow\res{k=-1}$.


% =====================================================
\seccio{10}{Factor comú i identitats notables}

\exf{1} \begin{tasks}(3)
  \task $\res{4x(x-2)}$
  \task $\res{3x\left(x^2+2x-3\right)}$
  \task $\res{5x^3(x-2)}$
  \task $\res{6x^2\left(2x^3-3x+1\right)}$
  \task $\res{x\left(x^2+x+1\right)}$
  \task $\res{2x^2\left(x^2-7\right)}$
\end{tasks}

\ex{1} \begin{tasks}(3)
  \task $\res{(x+6)(x-6)}$
  \task $\res{(x+4)^2}$
  \task $\res{(x-6)^2}$
  \task $\res{(3x+2)(3x-2)}$
  \task $\res{(2x+5)^2}$
  \task $\res{(5x-1)^2}$
  \task $\res{\left(x^2+3\right)\left(x^2-3\right)}$
  \task $\res{\left(x+\frac12\right)\left(x-\frac12\right)}$
  \task $\res{(4+x)(4-x)}$
\end{tasks}

\ex{2} \begin{tasks}(2)
  \task $2x\left(x^2-9\right)=\res{2x(x+3)(x-3)}$
  \task $3x\left(x^2+4x+4\right)=\res{3x(x+2)^2}$
  \task $x^2\left(x^2-1\right)=\res{x^2(x+1)(x-1)}$
  \task $5x\left(x^2-2x+1\right)=\res{5x(x-1)^2}$
  \task $x\left(x^4-16\right)=\res{x\left(x^2+4\right)(x+2)(x-2)}$
  \task $2x^2\left(x^2-25\right)=\res{2x^2(x+5)(x-5)}$
\end{tasks}

\ex{3} $83^2-17^2=(83+17)(83-17)=100\cdot66=\res{6600}$.


% =====================================================
\seccio{11}{Factorització de polinomis}

\exf{2} \begin{tasks}(2)
  \task $\res{(x-2)(x-3)}$
  \task $\res{(x+4)(x-3)}$
  \task $\res{2(x-1)(x-3)}$
  \task $\res{(x+1)(x-5)}$
\end{tasks}

\ex{2} \begin{tasks}(2)
  \task $\res{(x+2)(x-1)(x-3)}$
  \task $\res{(x+3)(x+2)(x-1)}$
  \task $\res{(x+2)(x+1)(x-3)}$
  \task $\res{(x+2)(x-2)(x-3)}$
\end{tasks}

\ex{3} \begin{tasks}(2)
  \task $\res{(x+2)(x+1)(x-1)(x-3)}$
  \task $\res{2(x+1)(x-1)(x-2)}$
  \task $\res{x^{2}(x+1)(x-3)}$
  \task $\res{(x+3)(x-2)^{2}}$
\end{tasks}

\ex{3} \begin{tasks}(1)
  \task $\res{(x-1)^{3}}$: $x=1$ (triple)
  \task $\res{(x+2)^{2}(x-2)^{2}}$: $x=-2$ i $x=2$ (dobles)
\end{tasks}


% =====================================================
\seccio{12}{Problemes}

\exf{2} $(x+2)^2-x^2=x^2+4x+4-x^2=\res{4x+4}$; \ per a $x=3$: $\res{16\ \text{cm}^2}$.

\ex{2} $V(x)=x(20-2x)(30-2x)=\res{4x^{3}-100x^{2}+600x}$; \ $V(5)=5\cdot10\cdot20=\res{1000\ \text{cm}^3}$.

\ex{3} $\left(x^2+5x+6\right):(x+2)=\res{x+3}$; \ $P(x)=2(x+2)+2(x+3)=\res{4x+10}$; \ $P(4)=\res{26}$.

\ex{3} $(x+3)^2-x^2=x^2+6x+9-x^2=\res{6x+9}$.


\end{document}
